\documentclass[../../../include/open-logic-section]{subfiles}

\begin{document}

\olfileid{sth}{cardinals}{cardsasords}
\olsection{اصلي عددونه د ترتيبي عددونو په توګه}

په حقيقت کښې به زموږ د اصلي عددونو نظريه په څرګنده توګه د ترتيبي
عددونو له نظريې کار اخلي۔ يعنې: اصلي عددونه به د ځانګړو ترتيبي
عددونو په توګه تعريف کړو۔ په تېره بيا، لاندې تعريف وړاندې کوو:

\begin{defn}\ollabel{defcardinalasordinal}
که $A$ ښه ترتيبېدے شي، نو $\card{A}$ تر ټولو کوچنے ترتيبي عدد
$\gamma$ دے چې $\cardeq{A}{\gamma}$ وي۔ د هر ترتيبي عدد $\gamma$
دپاره وايو چې $\gamma$ يو \emph{اصلي عدد} دے که او يوازې که
$\gamma = \card{\gamma}$ وي۔
\end{defn}

موږ اوس «$A$ ښه ترتيبېدے شي» وويل۔ لکه په رياضياتو کښې چې نږدې
تل وي، دلته د امکان عبارت يوازې د يوې وجودي ادعا لنډ بيان دے:
«$A$ ښه ترتيبېدے شي» دا وايي چې «داسې يوه اړيکه شته چې پر
$A$ ښه ترتيب جوړوي»۔

خو \olref{defcardinalasordinal} يوه ستونزه لري۔ غواړو چې \emph{هر}
سټ اندازه ولري، يعنې $\card{A}$ د هر~$A$ دپاره موجود وي۔ خو زموږ
تعريف په يوه شرط پيلېږي: «که $A$ ښه ترتيبېدے شي\ldots»۔ که
کوم سټ $A$ ښه ترتيب نه شي، نو دا تعريف د هغه دپاره $\card{A}$
څيز نه ټاکي۔

نو د \olref{defcardinalasordinal} د کارولو دپاره پکار ده چې هر سټ
ښه ترتيبېدے شي۔ خو له بده مرغه، په~$\ZF$ کښې دا تضمين نه لرو۔
له همدې امله، که \olref{defcardinalasordinal} کاروو، بايد يو نوے
اصل هم ور زيات کړو، لکه:
\begin{axiom}[ښه ترتيب]
هر سټ ښه ترتيبېدے شي۔
\end{axiom}
په \olref[choice][]{chap} کښې به وڅېړو چې د ښه ترتيب اصل د منلو
وړ دے که نه۔ خو له دې ځايه به يې کاروو۔ په دې اصل سره دا ثابتول
اسانه دي چې اصلي عددونه (د \olref{defcardinalasordinal} د تعريف
له مخې) موجود دي او ښه چلند کوي:

\begin{lem}\ollabel{lem:CardinalsExist}
د هر سټ $A$ دپاره:
\begin{enumerate}
	\item\ollabel{cardaexists} $\card{A}$ موجود او يوازينے دے؛
	\item\ollabel{cardaapprox}  $\cardeq{\card{A}}{A}$؛
	\item\ollabel{cardaidem}  $\card{A}$ اصلي عدد دے، يعنې
	$\card{A} = \card{\card{A}}$؛
\end{enumerate}
\end{lem}

\begin{proof}
$A$ وټاکئ۔ د ښه ترتيب له مخې $\tuple{A, R}$ يو ښه ترتيب
شته۔ د \olref[ordinals][ordtype]{thmOrdinalRepresentation} له مخې
$\tuple{A, R}$ له يوه يوازيني ترتيبي عدد $\beta$ سره هم‌شکل دے۔
نو $\cardeq{A}{\beta}$۔ د نامتناهيت هاخوا استقرا له مخې داسې يو
تر ټولو کوچنے ترتيبي عدد $\gamma$ شته چې $\cardeq{A}{\gamma}$
وي۔ نو $\card{A} = \gamma$، او \olref{cardaexists} او
\olref{cardaapprox} ثابت شول۔ د \olref{cardaidem} د ثبوت دپاره
وګورئ چې که $\delta \in \gamma$ وي، نو $\cardless{\delta}{A}$ د
$\gamma$ د ټاکنې له مخې؛ او ځکه چې هم‌شمېري يوه د همارزښت
اړيکه ده (\olref[sfr][siz][equ]{equinumerosityisequi})، نو
$\cardless{\delta}{\gamma}$ هم دے۔ نو $\gamma =
\card{\gamma}$۔
%So, for reductio, suppose that there is some ordinal $\delta \in \gamma$ such that $\cardeq{\gamma}{\delta}$. Then, $\cardeq{\cardeq{A}{\gamma}}{\delta}$ so that $\cardeq{A}{\delta}$ by \olref[sfr][set][equ]{equinumerosityisequi}, which contradicts the choice of $\gamma$ as the least ordinal such that $\cardeq{A}{\gamma}$. 
\end{proof}

راتلونکې پايله د کانتور اصل، او تر هغه زيات، تضمينوي۔ (په پام
کښې ونيسئ چې اصلي عددونه خپل ترتيب له ترتيبي عددونو اخلي، يعنې
$\cardfont{a} < \cardfont{b}$ هغه وخت او يوازې هغه وخت چې
$\cardfont{a} \in \cardfont{b}$ وي۔ د دې د بيان دپاره به د هغو
څيزونو دپاره چې اصلي عددونه يې ګڼو، فرکټور توري کاروو۔ دا عام
دود دے۔ يو بل عام دود د يوناني تورو کارول دي، ځکه چې اصلي عددونه
ترتيبي عددونه هم دي؛ خو د الفبا د منځ توري ټاکل کېږي، لکه
$\kappa, \lambda$۔):
\begin{lem}\ollabel{lem:CardinalsBehaveRight}
د هر ډول سټونو $A$ او $B$ دپاره:
\begin{align*}
	\cardeq{A}{B} &\text{ هغه وخت او يوازې هغه وخت چې } \card{A} = \card{B}\\
	\cardle{A}{B} &\text{ هغه وخت او يوازې هغه وخت چې } \card{A} \leq \card{B}\\
	\cardless{A}{B}&\text{ هغه وخت او يوازې هغه وخت چې } \card{A} < \card{B}
\end{align*}
\end{lem}

\begin{proof}
د دويمې ادعا له چپې خوا ښۍ خوا ته لورے ثابتوو (نور حالتونه ورته
دي او تمرين ته يې پرېږدو)۔ نو لاندې شکل وګورئ:
\begin{center}
	\begin{tikzpicture}
	\node (nodea) {$A$};
	\node[right = 6em of nodea] (nodeb) {$B$};
	\node[below = 2em of nodea] (nodecarda) {$\card{A}$};
	\node[below = 2em of nodeb] (nodecardb) {$\card{B}$};
	\draw[->] (nodea)--(nodeb);
	\draw[<->] (nodea)--(nodecarda);
	\draw[<->] (nodeb)--(nodecardb);
	\draw[->, dashed] (nodecarda)--(nodecardb);
\end{tikzpicture}
\end{center}
دواړه‌سري غشي !!{bijection}s ښيي چې د شتون تضمين يې
\olref{lem:CardinalsExist} کوي۔ له $\cardle{A}{B}$ څخه
!!a{injection} $A\to B$ شته۔ اوس د غشو په تعقيب له $\card{A}$
څخه $A$ ته، بيا $B$ ته، او بيا $\card{B}$ ته ځو؛ په دې توګه
!!a{injection} $\card{A} \to \card{B}$ ترلاسه کوو (مات غشے)۔
\end{proof}\noindent له \olref{lem:CardinalsBehaveRight} څخه د
شرودر--برنشتاين قضيه هم بيا ثابتولے شو۔ دا ادعا ده چې که
$\cardle{A}{B}$ او $\cardle{B}{A}$ وي، نو $\cardeq{A}{B}$۔
دا مو د \olref[sfr][siz][sb]{thm:schroder-bernstein} په توګه
بيان کړې وه، خو لومړے مو په \olref[sfr][infinite][card-sb]{sec}
کښې په يو څه هڅه ثابته کړه۔ اوس وګورئ:

\begin{proof}[د شرودر--برنشتاين بيا ثبوت]
که $\cardle{A}{B}$ او $\cardle{B}{A}$ وي، نو د
\olref{lem:CardinalsBehaveRight} له مخې $\card{A} \leq \card{B}$
او $\card{B} \leq \card{A}$۔ نو د درې‌ګوني وېش او
\olref{lem:CardinalsBehaveRight} له مخې $\card{A} = \card{B}$
او $\cardeq{A}{B}$۔
\end{proof}
\noindent
که څه هم دا ثبوت ډېر ساده دے، خو په ناڅرګنده توګه پر تعويض
(د \olref[ordinals][ordtype]{thmOrdinalRepresentation} د تضمين
دپاره) او پر ښه ترتيب (د \olref{lem:CardinalsBehaveRight} د
تضمين دپاره) دواړو تکيه کوي۔ پر خلاف يې، د
\olref[sfr][infinite][card-sb]{sec} ثبوت ډېر خپلواک و (په رښتيا
هم په~$\Zminus$ کښې ترسره کېدے شي)۔

\end{document}
